\section{Discussion}
\label{sec:discussion}

The strongest pattern in this run is a split between vulnerability signal and
execution stability. Gemini and Mistral variants occupy the higher-risk or
higher-hit region, while Claude and GPT variants show more stable execution with
lower hit rates.

\begin{figure}[t]
  \centering
  \includegraphics[width=\columnwidth]{frontier.pdf}
  \caption{Frontier view for the curated run.}
  \label{fig:frontier}
\end{figure}

From an engineering perspective, this reinforces a practical reporting rule:
aggregate risk, hit/refusal rates, and error rates should be interpreted
together rather than collapsed into one scalar.

A second methodological point is artifact discipline. The current paper is tied
strictly to archived run outputs under \path{reports/publication_run/} and
\path{trajectories/publication_run/}. This keeps claims auditable and limits
interpretive drift.

Finally, the engine package provides a path beyond broad campaign comparison
alone. Recon-led prioritization and confirmation reruns should improve
cost-efficiency and finding quality in future sweeps over the same curated
registry.
